\documentclass{article}
\usepackage{amsthm}
\usepackage{amsmath}
\usepackage{amssymb}
\newtheorem{thm}{Theorem}
\newtheorem{lem}[thm]{Lemma}

\begin{document}

\title{Proof of a Conjecture raised by Wu on OEIS Sequence A262814}
\author{Fung Cheok Yin}
\date{June 04th, 2026}

\maketitle

Definition of A262814: Numbers $k$ dividing every cyclic permutation of $k^k$. (As of June 04th, 2026).

To clearly understand the definition of the sequence, let us see Example from OEIS page: ``7 is a member as the six cyclic permutations of $7^7 = 823543$ are \{823543, 382354, 438235, 543823, 354382, 235438\} and these 6 integers are divisible by 7.''

Let $X$ be a non-negative integer and $N$ be its number of digits in base-10. We can define the cyclic permutation rigorously. For $X < 10$, $X$ has only one permutation, the digit itself. For $X \ge 10$, label its cyclic permutations by $X_0, X_1, \dots X_{N-1}$, which are generated by $X_{i+1} = (X_i - x_i + 10^N \cdot x_i)/10$ , for $i = 0 \dots N-2$ and $X = X_0$, $x_i = X_i \;\textrm{mod}\; 10$.


\begin{thm}
$10^n-1$ is a term of the sequence for all $n > 0$. (Conjecture raised by Chai Wah Wu on OEIS, Nov 03 2015.)
\end{thm}

\begin{lem}
Let $M = \;\textrm{Number of digits of}\; k^k$, where $k = 10^n-1$, then $M=nk$.
\end{lem}

\begin{proof}

The following proof follows from the proof of $99^{99}$ has 198 digits on Math StackExchange (https://math.stackexchange.com/a/2807815), by Barry Cipra. \\ In elementary real analysis, we have the well-known result $(1+\frac{1}{r})^r < e$ for all $r \in \mathbb{Z}^+$.

We can deduce, for all $r \in \mathbb{Z}^+$,

\[  (\frac{r}{r+1})^r > \frac{1}{e} \]

Then we have 
\begin{align*}
k^k & = \left( 10^n \left( \frac{10^n-1}{10^n} \right) \right)^k \\
& = 10^{nk} \left( \frac{10^n-1}{10^n} \right)^k \\
& = 10^{nk} \left( \frac{10^n-1}{10^n} \right)^k \\
& = 10^{nk} \left( \frac{10^n-1}{10^n} \right)^k \\
& = 10^{nk} \left( \frac{k}{k+1} \right)^k \\
& > \frac{10^{nk}}{e} \\
& > \frac{10^{nk}}{10} = 10^{nk-1}.
\end{align*}


On the other hand, 

\[ k^k < (10^n)^k = 10^{nk}. \]

Since $10^{nk} > k^k > 10^{nk-1}$ and the fact that $10^{nk-1}$ has $nk$ digits, we know the number of digits of $k^k$ is $nk$.

\end{proof}


\begin{lem}
Let $k = 10^n-1$ and $M = nk$. Assume a non-negative number $Y$ is divisible by $k$. Let $y$ be the unit digit of $Y$, the integer $Y'=(Y-y+10^M \cdot y)/10$ is divisible by $k$.
\end{lem}

\begin{proof}

Rewrite $Y'$ as $(Y+(10^M-1)y)/10$.

$ 10^M-1 = 10^{nk}-1 = (10^n-1)(10^{n(k-1)}+10^{n(k-2)}+\dots+1) = k(10^{n(k-1)}+10^{n(k-2)}+\dots+1)$, so $k|10^M-1$ and $k|(10^M-1)y$. From our assumption, we have $k|Y$. Combining these, we have $\;k\;|\;Y+(10^M-1)y$ . And $10 = 2 \cdot 5$, $\mathrm{gcd}(5, k) = 1$, $\mathrm{gcd}(2, k) = 1$, thus $\mathrm{gcd}(10, k) = 1$. Finally, from the basic rules of modular arithmetic we have $k|Y'$.

\end{proof}

\begin{proof}[Proof of Theorem 1]

It follows from the rigorous definition of cyclic permutations, Lemma 2 and Lemma 3.

\end{proof}

\end{document}
